\subsection{Compact Operators}

% **Role:** infinite-dimensional linear algebra powered by analysis.
%
% Compact maps, finite-rank approximation, Fredholm alternative, spectral theory of compact operators.

Compact operators carry bounded sets to relatively compact sets. Their spectral and solvability theory therefore resembles matrix theory inside infinite-dimensional spaces.

\begin{definition} \label{def:compact-and-finite-rank-operators} \hlt{Compact and Finite-Rank Operators}\\
Let $\mathcal{X},\mathcal{Y}$ be normed spaces and let $T:\mathcal{X}\to\mathcal{Y}$ be linear.
\begin{enumerate}[label=(\roman*)]
\setlength{\itemsep}{0pt}
\item \hlt{Compact}: $\overline{T(B)}$ is compact in $\mathcal{Y}$ $\forall B\subseteq\mathcal{X}$ bounded.
\item \hlt{Finite rank}: $\dim\operatorname{range}(T)<\infty$.
\end{enumerate}
\end{definition}

\begin{proposition} \label{prop:compact-operator-total-bounded-sequential-characterizations} \hlt{Compactness Tests on Hilbert Spaces}\\
Let $\mathcal{H}$ be Hilbert, let $K\in\LL(\mathcal{H})$, and set $B\defeq\overline{B}_{\mathcal{H}}$. The following are equivalent:
\begin{enumerate}[label=(\roman*)]
\setlength{\itemsep}{0pt}
\item $K$ is compact.
\item $K(B)$ is totally bounded.
\item Every bounded $(h_n)\subseteq\mathcal{H}$ has a subsequence $(h_{n_j})$ such that $Kh_{n_j}$ converges in norm.
\end{enumerate}
\begin{proof}
For bounded $E\subseteq\mathcal{H}$, choose $r>0$ with $E\subseteq rB$; then $K(E)\subseteq K(rB)=rK(B)$.
\begin{enumerate}[label=(\roman*)]
\setlength{\itemsep}{0pt}
\item[$(i)$]$\Rightarrow(ii)$ $K$ compact $\Rightarrow \overline{K(B)}$ compact $\Rightarrow \overline{K(B)}$ totally bounded $\Rightarrow K(B)$ totally bounded.
\item[$(ii)$]$\Rightarrow(iii)$ $(h_n)$ bounded $\Rightarrow \exists r>0$, $\{h_n:n\geq1\}\subseteq rB$. Then $(Kh_n)\subseteq K(rB)=rK(B)$ with $rK(B)$ totally bounded, so $\exists(n_j)$ such that $(Kh_{n_j})$ is Cauchy; $\mathcal{H}$ complete $\Rightarrow \exists y\in\mathcal{H}$, $\norm{Kh_{n_j}-y}\to0$.
\item[$(iii)$]$\Rightarrow(i)$ Let $E\subseteq\mathcal{H}$ be bounded. If $K(E)$ is not totally bounded, then $\exists\varepsilon>0$ and $(x_n)\subseteq E$ such that $n\neq m\Rightarrow\norm{Kx_n-Kx_m}\geq\varepsilon$. By (iii), $\exists(n_j)$ and $y\in\mathcal{H}$ with $\norm{Kx_{n_j}-y}\to0$, contradiction. Hence $K(E)$ is totally bounded, so $\overline{K(E)}$ is totally bounded and complete; Theorem \ref{thm:compact-complete-totally-bounded} gives $\overline{K(E)}$ compact.
\end{enumerate}
\hlt{Structural Commentary}: $E\subseteq rB\Rightarrow K(E)\subseteq rK(B)$; $K(B)$ totally bounded $\Rightarrow (h_n)\subseteq rB\Rightarrow\exists(n_j):(Kh_{n_j})$ Cauchy.\\
\hlt{Key Insight}: $K$ compact iff $K$ sends bounded sequences to norm-precompact sequences.
\end{proof}
\end{proposition}

\begin{proposition} \label{prop:finite-rank-operators-compact} \hlt{Finite Rank Operators Are Compact}\\
Let $\mathcal{X},\mathcal{Y}$ be normed, $T\in\LL(\mathcal{X},\mathcal{Y})$,$R\defeq\operatorname{range}(T)$. If $\dim R<+\infty$, then $T$ is compact.
\begin{proof} Let $B\subseteq\mathcal{X}$ be bounded. Then $T(B)\subseteq R$ is bounded, so $\overline{T(B)}^{\,R}$ is compact by Theorem \ref{thm:finite-dimensional-normed-spaces}.\\
Since $R\leq\mathcal{Y}$ is finite-dimensional, $R$ is closed in $\mathcal{Y}$; hence $\overline{T(B)}^{\,\mathcal{Y}}=\overline{T(B)}^{\,R}$ is compact. Thus $T$ is compact.\\
\hlt{Structural Commentary}: $T(B)\subseteq R$, $\dim R<+\infty$, and bounded subsets of $R$ have compact closure.\\
\hlt{Key Insight}: Finite-dimensional range makes bounded images relatively compact.
\end{proof}
\end{proposition}

\begin{theorem} \label{thm:compact-operator-closed-under-norm-limits} \hlt{Norm Limits of Compact Operators}\\
Let $\mathcal{X}$ be normed, $\mathcal{Y}$ Banach, $T_n,T\in\LL(\mathcal{X},\mathcal{Y})$. If each $T_n$ compact and $\normop{T_n-T}\to0$, then $T$ is compact.
\begin{proof} Let $B\subseteq\mathcal{X}$ be bounded, $M>0$, $\norm{x}\leq M$ for $x\in B$. For $\varepsilon>0$, choose $N$ with $\normop{T-T_N}<\frac{\varepsilon}{2M}$.\\
Since $\overline{T_N(B)}$ is compact, choose $y_1,\ldots,y_r\in\mathcal{Y}$ with $T_N(B)\subseteq\bigcup_{j=1}^r B_{\mathcal{Y}}(y_j,\frac{\varepsilon}{2})$.\\
Then $T(B)\subseteq\bigcup_{j=1}^r B_{\mathcal{Y}}(y_j,\varepsilon)$, so $T(B)$ is totally bounded. Since $\mathcal{Y}$ is complete, $\overline{T(B)}$ is compact; $T$ is compact.
\hlt{Structural Commentary}: $\normop{T-T_N}$ transfers finite nets from $T_N(B)$ to $T(B)$ for bounded $B$.\\
\hlt{Key Insight}: Norm limits preserve relative compactness of bounded images.
\end{proof}
\end{theorem}

\begin{proposition} \label{prop:compact-operators-linear-ideal} \hlt{Compact Operators Form an Operator Ideal}\\
Let $\mathcal{W},\mathcal{X},\mathcal{Y},\mathcal{Z}$ be normed spaces.
\begin{enumerate}[label=(\roman*)]
\setlength{\itemsep}{0pt}
\item If $K_1,\ldots,K_m:\mathcal{X}\to\mathcal{Y}$ are compact, then $\sum\limits_{j=1}^m K_j$ is compact.
\item If $K:\mathcal{X}\to\mathcal{Y}$ is compact, $A\in\LL(\mathcal{Y},\mathcal{Z})$, and $B\in\LL(\mathcal{W},\mathcal{X})$, then $AKB:\mathcal{W}\to\mathcal{Z}$ is compact.
\end{enumerate}
\begin{proof}
\begin{enumerate}[label=(\roman*)]
\setlength{\itemsep}{0pt}
\item Let $E\subseteq\mathcal{X}$ be bounded, $C_j\defeq\overline{K_j(E)}$. Then each $C_j$ is compact, and $\sum\limits_{j=1}^m K_j(E)\subseteq \sigma(C_1\times\cdots\times C_m)$ for the continuous map $\sigma(y_1,\ldots,y_m)\defeq\sum\limits_{j=1}^m y_j$. Thus $\sum\limits_{j=1}^m K_j(E)$ is relatively compact.
\item Let $E\subseteq\mathcal{W}$ be bounded. Then $B(E)$ is bounded, $C\defeq\overline{K(B(E))}$ is compact, and $AKB(E)\subseteq A(C)$.\\
Since $A(C)$ is compact and closed in $\mathcal{Z}$, $\overline{AKB(E)}\subseteq A(C)$ is compact.
\end{enumerate}
\hlt{Structural Commentary}: Products of compact image closures control sums; bounded precomposition and continuous postcomposition control $AKB$.\\
\hlt{Key Insight}: Compact operators form a two-sided ideal in $\LL$.
\end{proof}
\end{proposition}

\begin{example} \label{ex:diagonal-compact-operator-orthonormal-basis} \hlt{Diagonal Compact Operators}\\
Let $(e_k)_{k\geq1}$ be a Hilbert basis of separable $\mathcal{H}$ and let $\lambda=(\lambda_k)\in\ell^\infty$ with $\lambda_k\to0$.\\
For $h\in\mathcal{H}$, set $Th\defeq\sum\limits_{k=1}^{\infty}\lambda_k\innerprod{h,e_k}e_k$ and $T_Nh\defeq\sum\limits_{k=1}^{N}\lambda_k\innerprod{h,e_k}e_k$.\\
By Parseval: $\sum\limits_k\abs{\lambda_k\innerprod{h,e_k}}^2\leq\norm{\lambda}_\infty^2\norm{h}^2\Rightarrow T\in\LL(\mathcal{H})$.\\
For $N\geq1$, $\operatorname{range}T_N\subseteq\operatorname{span}\{e_1,\ldots,e_N\}\Rightarrow\operatorname{rank}T_N\leq N\Rightarrow T_N$ compact by Proposition \ref{prop:finite-rank-operators-compact}.\\
Also $\norm{(T-T_N)h}^2=\sum\limits_{k>N}\abs{\lambda_k}^2\abs{\innerprod{h,e_k}}^2\leq\sup\limits_{k>N}\abs{\lambda_k}^2\norm{h}^2\Rightarrow\normop{T-T_N}\leq\sup\limits_{k>N}\abs{\lambda_k}\to0$.\\
Hence $T_N\to T$ in $\normop{\cdot}$ and $T$ is compact by Theorem \ref{thm:compact-operator-closed-under-norm-limits}.
\end{example}

\begin{theorem} \label{thm:compact-operator-weak-to-strong-characterization} \hlt{Weak-to-Strong Characterization}\\
Let $\mathcal{H}$ be Hilbert, $K\in\LL(\mathcal{H})$. Then $K$ is compact iff $h_n\rightharpoonup h$ implies $\norm{Kh_n-Kh}\to0$ $\forall (h_n)\subseteq\mathcal{H}$ and $h\in\mathcal{H}$.
\begin{proof}
\begin{enumerate}[label=\roman*.]
\setlength{\itemsep}{0pt}
\item[$(\Rightarrow)$] Let $h_n\rightharpoonup h$. By Proposition \ref{prop:weakly-convergent-sequences-bounded}(i), $(h_n)$ bounded.\\
For every subsequence $(n_j)$, By Proposition \ref{prop:compact-operator-total-bounded-sequential-characterizations}(iii), $\exists(j_\ell)$ and $y\in\mathcal{H}$ with $\norm{Kh_{n_{j_\ell}}-y}\to0$.\\
For $u\in\mathcal{H}$, $\innerprod{y,u}=\lim\limits_{\ell}\innerprod{Kh_{n_{j_\ell}},u}=\lim\limits_{\ell}\innerprod{h_{n_{j_\ell}},K^*u}=\innerprod{h,K^*u}=\innerprod{Kh,u}$, so $y=Kh$.\\
If $\norm{Kh_n-Kh}\not\to0$, then $\exists\varepsilon>0$, $(n_j)$ with $\norm{Kh_{n_j}-Kh}\geq\varepsilon$, contradicting $\exists(j_\ell):\norm{Kh_{n_{j_\ell}}-Kh}\to0$.
\item[$(\Leftarrow)$] Let $(h_n)$ be bounded. By Theorem \ref{thm:hilbert-weak-subsequential-compactness}, $\exists(n_j)$, $h\in\mathcal{H}$ with $h_{n_j}\rightharpoonup h$. By hypothesis, $\norm{Kh_{n_j}-Kh}\to0$.\\
Hence $\forall(h_n)$ bounded, $\exists(n_j)$ and $y\in\mathcal{H}$ with $\norm{Kh_{n_j}-y}\to0$; By Proposition \ref{prop:compact-operator-total-bounded-sequential-characterizations}(iii), $K$ compact.
\end{enumerate}
\hlt{Structural Commentary}: $K$ compact, $h_n\rightharpoonup h\Rightarrow Kh_{n_{j_\ell}}\to y=Kh$; $(h_n)$ bounded $\Rightarrow h_{n_j}\rightharpoonup h\Rightarrow Kh_{n_j}\to Kh$.\\
\hlt{Key Insight}: $K$ compact iff $h_n\rightharpoonup h$ forces $Kh_n\to Kh$ in norm.
\end{proof}
\end{theorem}

\begin{definition} \label{def:orthogonal-projection-sequence-hilbert-basis} \hlt{Orthogonal Projection Sequence}\\
Let $(e_k)_{k\geq1}$ be a Hilbert basis for a separable Hilbert space $\mathcal{H}$.\\
The induced \hlt{orthogonal projection sequence} $(P_n)_{n\geq1}$ is defined by:
\begin{enumerate}[label=(\roman*)]
\setlength{\itemsep}{0pt}
\item $P_nh\defeq\sum\limits_{k=1}^{n}\innerprod{h,e_k}e_k$ for $h\in\mathcal{H}$.
\item $\operatorname{range}(P_n)=\operatorname{span}\{e_1,\ldots,e_n\}$ and $\normop{P_n}=1$.
\item $P_nh\to h$ for every $h\in\mathcal{H}$.
\end{enumerate}
\end{definition}

\begin{proposition} \label{prop:compact-operators-finite-rank-projection-approximation} \hlt{Finite-Rank Approximation by Orthogonal Projections}\\
Let $(P_n)$ be the orthogonal projection sequence induced by a Hilbert basis of a separable Hilbert space $\mathcal{H}$.\\
For $T\in\LL(\mathcal{H})$, $T$ is compact iff $\normop{P_nT-T}\to0$.
\begin{proof}
\begin{enumerate}[label=\roman*.]
\setlength{\itemsep}{0pt}
\item[$(\Rightarrow)$] $B\defeq\overline{B}_{\mathcal{H}}$, $C\defeq\overline{T(B)}$; $T$ compact $\Rightarrow C$ compact.\\
$\varepsilon>0\Rightarrow\exists y_1,\ldots,y_r\in C,N\in\mathbb{N}:C\subseteq\bigcup\limits_{j=1}^r B_{\mathcal{H}}(y_j,\frac{\varepsilon}{3})$ and $n\geq N\Rightarrow\max\limits_{1\leq j\leq r}\norm{(P_n-I)y_j}<\frac{\varepsilon}{3}$.\\
Thus $n\geq N$, $x\in B\Rightarrow\exists j:\norm{Tx-y_j}<\frac{\varepsilon}{3}$ and $\norm{(P_nT-T)x}\leq\norm{Tx-y_j}+\norm{(P_n-I)y_j}+\norm{y_j-Tx}<\varepsilon$.\\
Hence $\normop{P_nT-T}=\sup\limits_{x\in B}\norm{(P_nT-T)x}\to0$.
\item[$(\Leftarrow)$] $\operatorname{range}(P_nT)\subseteq\operatorname{range}(P_n)$, $\dim\operatorname{range}(P_n)<\infty\Rightarrow P_nT$ compact by Proposition \ref{prop:finite-rank-operators-compact}.\\
Also $P_nT\to T$ in $\normop{\cdot}$; Theorem \ref{thm:compact-operator-closed-under-norm-limits} $\Rightarrow T$ compact.
\end{enumerate}
\hlt{Structural Commentary}: Compact $C=\overline{T(B)}$ upgrades $P_ny\to y$ to uniform convergence on $C$.\\
\hlt{Key Insight}: $T$ compact iff finite-rank projections approximate $T$ in operator norm.
\end{proof}
\end{proposition}

\begin{corollary} \label{cor:compact-operators-finite-rank-dense-separable-hilbert} \hlt{Finite-Rank Operators Are Dense in Compact Operators}\\
Let $\mathcal{H}$ be separable, and let $\mathcal{K}(\mathcal{H})$ and $\mathcal{F}(\mathcal{H})$ denote the compact and finite-rank operators in $\LL(\mathcal{H})$.
\begin{enumerate}[label=(\roman*)]
\setlength{\itemsep}{0pt}
\item $\mathcal{K}(\mathcal{H})$ is a norm-closed linear subspace of $\LL(\mathcal{H})$.
\item $\overline{\mathcal{F}(\mathcal{H})}^{\,\normop{\cdot}}=\mathcal{K}(\mathcal{H})$.
\end{enumerate}
\begin{proof}
\begin{enumerate}[label=(\roman*)]
\setlength{\itemsep}{0pt}
\item Proposition \ref{prop:compact-operators-linear-ideal}: $K_1,K_2\in\mathcal{K}(\mathcal{H})$, $\alpha,\beta\in\mathbb{F}\Rightarrow\alpha K_1+\beta K_2\in\mathcal{K}(\mathcal{H})$.\\
Theorem \ref{thm:compact-operator-closed-under-norm-limits}: $K_n\in\mathcal{K}(\mathcal{H})$, $\normop{K_n-K}\to0\Rightarrow K\in\mathcal{K}(\mathcal{H})$.
\item[$(ii)$]$(\Rightarrow)$ $\mathcal{F}(\mathcal{H})\subseteq\mathcal{K}(\mathcal{H})$, (i) $\Rightarrow\overline{\mathcal{F}(\mathcal{H})}^{\,\normop{\cdot}}\subseteq\mathcal{K}(\mathcal{H})$.\\
\item[$(ii)$]$(\Leftarrow)$ $\dim\mathcal{H}<+\infty\Rightarrow\mathcal{F}(\mathcal{H})=\LL(\mathcal{H})\supseteq\mathcal{K}(\mathcal{H})$. If $\dim\mathcal{H}=+\infty$,\\
take $(P_n)$; $T\in\mathcal{K}(\mathcal{H})\Rightarrow P_nT\in\mathcal{F}(\mathcal{H})$, $\normop{P_nT-T}\to0\Rightarrow T\in\overline{\mathcal{F}(\mathcal{H})}^{\,\normop{\cdot}}$ by Proposition \ref{prop:compact-operators-finite-rank-projection-approximation}.
\end{enumerate}
\hlt{Structural Commentary}: $\mathcal{K}(\mathcal{H})$ is norm closed, and $P_nT\in\mathcal{F}(\mathcal{H})$ with $P_nT\to T$ for compact $T$.\\
\hlt{Key Insight}: On separable $\mathcal{H}$, $\mathcal{K}(\mathcal{H})=\overline{\mathcal{F}(\mathcal{H})}^{\,\normop{\cdot}}$.
\end{proof}
\end{corollary}

\begin{theorem} \label{thm:schauder-compact-adjoint} \hlt{Schauder's Theorem}\\
Let $\mathcal{H}$ be Hilbert. If $K\in\LL(\mathcal{H})$ is compact, then $K^*$ is compact.
\begin{proof} ByTheorem \ref{thm:hilbert-adjoint-existence-bounded},$K^*\in\LL(\mathcal{H})$. By Theorem \ref{thm:compact-operator-weak-to-strong-characterization}, it suffices to prove $h_n\rightharpoonup h\Rightarrow\norm{K^*(h_n-h)}\to0$.\\
Put $u_n\defeq h_n-h$. Then $u_n\rightharpoonup0$, $M\defeq\sup_n\norm{u_n}<\infty$, and $\forall v\in\mathcal{H}$, $\innerprod{K^*u_n,v}=\innerprod{u_n,Kv}\to0$, so $K^*u_n\rightharpoonup0$.\\
Compactness of $K$ gives $\norm{KK^*u_n}\to0$, hence $0\leq\norm{K^*u_n}^2=\innerprod{KK^*u_n,u_n}\leq M\norm{KK^*u_n}\to0$.\\
\hlt{Structural Commentary}: $u_n\rightharpoonup0\Rightarrow K^*u_n\rightharpoonup0$, $K$ compact $\Rightarrow KK^*u_n\to0$; $\norm{K^*u_n}^2\leq M\norm{KK^*u_n}$.\\
\hlt{Key Insight}: $\norm{K^*u_n}^2=\innerprod{KK^*u_n,u_n}$.
\end{proof}
\end{theorem}

\begin{proposition} \label{prop:compact-identity-perturbation-injective-surjective} \hlt{Injective Compact Perturbations of the Identity Are Onto}\\
Let $\mathcal{H}$ be Hilbert, let $K\in\LL(\mathcal{H})$ be compact, and set $A\defeq I+K$. If $\ker A=\{0\}$, then $A(\mathcal{H})=\mathcal{H}$.
\begin{proof}
\begin{enumerate}[label=(\roman*)]
\setlength{\itemsep}{0pt}
\item If $\inf_{\norm{h}=1}\norm{Ah}=0$, choose $\norm{h_n}=1$, $Ah_n\to0$. By Proposition \ref{prop:compact-operator-total-bounded-sequential-characterizations}(iii), $Kh_{n_j}\to y$,\\
hence $h_{n_j}=Ah_{n_j}-Kh_{n_j}\to-y\defeq h$, $\norm{h}=1$, and $Ah=0$, contradiction.\\
Thus $\exists c>0:\norm{Ah}\geq c\norm{h}$ $\forall h\in\mathcal{H}$.
\item Fix $g\in\mathcal{H}$ and put $M\defeq\operatorname{cspan}\{K^m g:m\geq0\}$. Then $g\in M$, $K(M)\subseteq M$, and $M$ is separable.\\
If $\dim M<+\infty$, then $A|_M:M\to M$ injective $\Rightarrow A|_M$ surjective $\Rightarrow g\in A(\mathcal{H})$.
\item Assume $\dim M=+\infty$. Let $(P_n)$ be projection sequence of Hilbert basis of $M$,\\
$H_n\defeq P_n(M)$, $\dim H_n<+\infty$. By Proposition \ref{prop:compact-operators-finite-rank-projection-approximation}, $\normop{P_nK-K}\to0$. Let $N$,\\
$n\geq N\Rightarrow\normop{P_nK-K}<\frac{c}{2}$. Then $n\geq N$, $h\in H_n\Rightarrow\norm{(I+P_nK)h}\geq\norm{Ah}-\norm{(P_nK-K)h}\geq\frac{c}{2}\norm{h}$, so $I+P_nK:H_n\to H_n$ injective $\Rightarrow$ surjective.
\item Choose $u_n\in H_n$ with $(I+P_nK)u_n=P_ng$. Then $\norm{u_n}\leq\frac{2}{c}\norm{P_ng}\leq\frac{2}{c}\norm{g}$.\\
By Theorem \ref{thm:hilbert-weak-subsequential-compactness}, $u_{n_j}\rightharpoonup u$ for some subsequence.\\
By Theorem \ref{thm:compact-operator-weak-to-strong-characterization}, $Ku_{n_j}\to Ku$, $\norm{P_{n_j}Ku_{n_j}-Ku_{n_j}}\leq\normop{P_{n_j}K-K}\norm{u_{n_j}}\to0$.\\
Since $P_{n_j}g\to g$, $u_{n_j}=P_{n_j}g-P_{n_j}Ku_{n_j}\to g-Ku$; by weak uniqueness $u=g-Ku$, hence $Au=g$.
\end{enumerate}
Since $g\in\mathcal{H}$ was arbitrary, $A(\mathcal{H})=\mathcal{H}$.\\
\hlt{Structural Commentary}: $\norm{Ah}\geq c\norm{h}$ and $(I+P_nK)u_n=P_ng\Rightarrow Au=g$.\\
\hlt{Key Insight}: Compactness removes approximate kernels and passes projected solutions to a true solution.
\end{proof}
\end{proposition}

\begin{definition} \label{def:fredholm-operator-index} \hlt{Fredholm Operator and Index}\\
Let $\mathcal{H}$ be a Hilbert space and let $T\in\LL(\mathcal{H})$.
\begin{enumerate}[label=(\roman*)]
\setlength{\itemsep}{0pt}
\item \hlt{Fredholm}: $\operatorname{range}(T)$ is closed, $\dim\ker T<\infty$, and $\operatorname{codim}\operatorname{range}(T)<\infty$.
\item \hlt{Index}: $\operatorname{ind}T\defeq\dim\ker T-\operatorname{codim}\operatorname{range}(T)$ when $T$ is Fredholm.
\end{enumerate}
\end{definition}

\begin{theorem} \label{thm:riesz-schauder-compact-perturbation-identity} \hlt{Riesz--Schauder for Compact Perturbations of the Identity}\\
Let $\mathcal{H}$ be Hilbert, let $K\in\LL(\mathcal{H})$ be compact, and set $A\defeq I+K$. Then:
\begin{enumerate}[label=(\roman*)]
\setlength{\itemsep}{0pt}
\item $\dim\ker A<\infty$.
\item $\operatorname{range}(A)$ is closed.
\item $\operatorname{codim}\operatorname{range}(A)=\dim\ker(A^*)=\dim\ker A<\infty$.
\item $A$ is Fredholm with $\operatorname{ind}A=0$.
\item $A$ is injective iff it is surjective.
\end{enumerate}
\begin{proof}
\begin{enumerate}[label=(\roman*)]
\setlength{\itemsep}{0pt}
\item $\dim\ker A=+\infty\Rightarrow\exists(u_n)\subseteq\ker A$ with $\innerprod{u_n,u_m}=\delta_{nm}$.\\
Then $Ku_n=-u_n$ and $n\neq m\Rightarrow\norm{Ku_n-Ku_m}=\sqrt2$, contradicting Proposition \ref{prop:compact-operator-total-bounded-sequential-characterizations}(iii).
\item Let $H_0\defeq(\ker A)^\perp$. If $\inf\{\norm{Ah}:h\in H_0,\norm{h}=1\}=0$, choose $h_n\in H_0$, $\norm{h_n}=1$, $Ah_n\to0$.\\
By Proposition \ref{prop:compact-operator-total-bounded-sequential-characterizations}(iii), $Kh_{n_j}\to y$, so $h_{n_j}=Ah_{n_j}-Kh_{n_j}\to-y\defeq h\in H_0$, $Ah=0$.\\
Thus $h\in\ker A\cap H_0=\{0\}$, contradiction. Hence $\exists c>0:\norm{Ah}\geq c\norm{h}$ on $H_0$.\\
Also $A(H_0)=A(\mathcal{H})$. If $Ah_n\to y$ with $h_n\in H_0$, then $\norm{h_n-h_m}\leq c^{-1}\norm{Ah_n-Ah_m}$.\\
Thus $h_n\to h\in H_0$ and $Ah=y$; hence $\operatorname{range}(A)$ is closed.
\item By Theorem \ref{thm:schauder-compact-adjoint}, $K^*$ compact, so (i)--(ii) apply to $A^*=I+K^*$.\\
By Proposition \ref{prop:closed-range-splitting-coercive-invertibility}, $\mathcal{H}=\operatorname{range}(A)\odot\ker(A^*)$, $\operatorname{codim}\operatorname{range}(A)=\dim\ker(A^*)$.\\
If $\dim\ker A<\dim\ker(A^*)$, choose injective $L:\ker A\to\ker(A^*)$ with $L(\ker A)\neq\ker(A^*)$, $F\defeq LP_{\ker A}$.\\
Then $\operatorname{rank}F<+\infty$, $K+F$ compact. For $x=z+n$, $z\in H_0$, $n\in\ker A$,\\
$(A+F)x=Az+Ln\in\operatorname{range}(A)\odot L(\ker A)$; injectivity of $L$ gives $\ker(A+F)=\{0\}$ and\\
$\operatorname{range}(A+F)=\operatorname{range}(A)\odot L(\ker A)\neq\mathcal{H}$, contradicting Proposition \ref{prop:compact-identity-perturbation-injective-surjective}. Thus $\dim\ker A\geq\dim\ker(A^*)$.\\
Apply the same argument to $A^*$ and use $(K^*)^*=K$ to get $\dim\ker A=\dim\ker(A^*)$.
\item (i)--(iii) give $\operatorname{range}(A)$ closed, $\dim\ker A<+\infty$, $\operatorname{codim}\operatorname{range}(A)<+\infty$.\\
Also $\dim\ker A=\operatorname{codim}\operatorname{range}(A)$; hence $A$ is Fredholm and $\operatorname{ind}A=0$.
\item (iii) gives $\ker A=\{0\}\Longleftrightarrow\operatorname{codim}\operatorname{range}(A)=0\Longleftrightarrow\operatorname{range}(A)=\mathcal{H}$.
\end{enumerate}
\hlt{Structural Commentary}: Compactness $\Rightarrow\dim\ker A<\infty$, $\operatorname{range}(A)$ closed; $F=LP_{\ker A}\Rightarrow\dim\ker A=\dim\ker(A^*)$.\\
\hlt{Key Insight}: For $A=I+K$ compactly perturbed, injectivity and surjectivity have the same finite defect.
\end{proof}
\end{theorem}

\begin{corollary} \label{cor:fredholm-alternative-compact-operator} \hlt{Fredholm Alternative}\\
Let $\mathcal{H}$ be Hilbert, let $K\in\LL(\mathcal{H})$ be compact, let $\lambda\in\mathbb{F}\setminus\{0\}$, and set $T_\lambda\defeq\lambda I-K$.
\begin{enumerate}[label=(\roman*)]
\setlength{\itemsep}{0pt}
\item $\operatorname{range}(T_\lambda)=\mathcal{H}$ iff $\ker T_\lambda=\{0\}$.
\item If the conditions in (i) hold, then $T_\lambda^{-1}\in\LL(\mathcal{H})$.
\end{enumerate}
\begin{proof} Let $\widetilde K\defeq-\lambda^{-1}K$; $\widetilde K$ compact, $T_\lambda=\lambda(I+\widetilde K)$, $\ker T_\lambda=\ker(I+\widetilde K)$, $\operatorname{range}(T_\lambda)=\mathcal{H}$ iff $\operatorname{range}(I+\widetilde K)=\mathcal{H}$.
\begin{enumerate}[label=(\roman*)]
\setlength{\itemsep}{0pt}
\item[$(i)$] ($\Rightarrow$) $\operatorname{range}(T_\lambda)=\mathcal{H}\Rightarrow\operatorname{range}(I+\widetilde K)=\mathcal{H}\Rightarrow\ker(I+\widetilde K)=\{0\}$ by Theorem \ref{thm:riesz-schauder-compact-perturbation-identity}(v), hence $\ker T_\lambda=\{0\}$.
\item[$(i)$] ($\Leftarrow$) $\ker T_\lambda=\{0\}\Rightarrow\ker(I+\widetilde K)=\{0\}\Rightarrow\operatorname{range}(I+\widetilde K)=\mathcal{H}$ by Theorem \ref{thm:riesz-schauder-compact-perturbation-identity}(v), hence $\operatorname{range}(T_\lambda)=\mathcal{H}$.
\item[$(ii)$] (i) $\Rightarrow T_\lambda$ bijective; Corollary \ref{cor:banach-isomorphism-theorem} $\Rightarrow T_\lambda^{-1}\in\LL(\mathcal{H})$.
\end{enumerate}
\hlt{Structural Commentary}: $T_\lambda=\lambda(I+\widetilde K)$ with $\widetilde K$ compact; bijective $T_\lambda$ has bounded inverse.\\
\hlt{Key Insight}: $\forall f$, $T_\lambda u=f$ is solvable iff $\ker T_\lambda=\{0\}$.
\end{proof}
\end{corollary}

\begin{theorem} \label{thm:fredholm-index-zero-invertible-plus-compact} \hlt{Fredholm Index Zero as an Invertible Plus Compact Operator}\\
Let $\mathcal{H}$ be Hilbert and let $T\in\LL(\mathcal{H})$. The following are equivalent:
\begin{enumerate}[label=(\roman*)]
\setlength{\itemsep}{0pt}
\item $T$ is Fredholm and $\operatorname{ind}T=0$.
\item $\exists S,K\in\LL(\mathcal{H})$ such that $T=S+K$, $S$ is invertible, and $K$ is compact.
\end{enumerate}
\begin{proof}
\begin{enumerate}[label=(\roman*)]
\setlength{\itemsep}{0pt}
\item[$(i)$]$\Rightarrow(ii)$ $N\defeq\ker T$, $R\defeq\operatorname{range}(T)$, $Z\defeq\ker(T^*)$.\\
By Proposition \ref{prop:closed-range-splitting-coercive-invertibility}, $\mathcal{H}=R\odot Z$ and $\dim Z=\operatorname{codim}R$. $\operatorname{ind}T=0\Rightarrow\dim N=\dim Z<+\infty$.\\
$A:N\to Z$ isomorphism, $P_N$ orthogonal projection, $K\defeq AP_N$, $S\defeq T-K$.\\
Then $\operatorname{rank}K<+\infty\Rightarrow K$ compact and $T=S+K$. For $h=v+u$, $v\in N^\perp$, $u\in N$, $Sh=Tv-Au\in R\odot Z$ and $Sh=0\Rightarrow Tv=0$, $Au=0\Rightarrow v=0$, $u=0$; hence $\ker S=\{0\}$.\\
Also $T(N^\perp)=R$, $A(N)=Z$, and $y=r+z\in R\odot Z\Rightarrow\exists v\in N^\perp,u\in N:Tv=r,Au=-z,S(v+u)=y$; hence $\operatorname{range}S=\mathcal{H}$ and $S^{-1}\in\LL(\mathcal{H})$ by Corollary \ref{cor:banach-isomorphism-theorem}.
\item[$(ii)$]$\Rightarrow(i)$ $C\defeq S^{-1}K$, $A\defeq I+C$. By Proposition \ref{prop:compact-operators-linear-ideal}(ii), $C$ is compact and $T=SA$.\\
By Theorem \ref{thm:riesz-schauder-compact-perturbation-identity}(iv), $A$ is Fredholm, $\operatorname{ind}A=0$. $S$ isomorphism $\Rightarrow \ker T=\ker A$, $\operatorname{range}T=S(\operatorname{range}A)$, $\operatorname{codim}\operatorname{range}T=\operatorname{codim}\operatorname{range}A$. Hence $T$ is Fredholm and $\operatorname{ind}T=0$.
\end{enumerate}
\hlt{Structural Commentary}: $K=AP_{\ker T}$ makes $S=T-K$ invertible; $T=S(I+C)$ with $C$ compact.\\
\hlt{Key Insight}: $\operatorname{ind}T=0$ iff $T$ is invertible modulo compact operators.
\end{proof}
\end{theorem}

\begin{corollary} \label{cor:fredholm-index-zero-closed-under-composition} \hlt{Index-Zero Fredholm Operators Are Closed under Composition}\\
Let $\mathcal{H}$ be Hilbert, $S,T\in\LL(\mathcal{H})$ be Fredholm with $\operatorname{ind}S=\operatorname{ind}T=0$. Then $ST$ is Fredholm and $\operatorname{ind}(ST)=0$.
\begin{proof} By Theorem \ref{thm:fredholm-index-zero-invertible-plus-compact}, $S=S_0+K_S$, $T=T_0+K_T$, with $S_0,T_0$ invertible and $K_S,K_T$ compact.\\
Let $U\defeq S_0T_0$, $C\defeq S_0K_T+K_ST_0+K_SK_T$. Then $ST=U+C$, $U$ invertible, by Proposition \ref{prop:compact-operators-linear-ideal}, $C$ compact.\\
Theorem \ref{thm:fredholm-index-zero-invertible-plus-compact} $\Rightarrow ST$ Fredholm and $\operatorname{ind}(ST)=0$.\\
\hlt{Structural Commentary}: $(S_0+K_S)(T_0+K_T)=S_0T_0+C$, $C\in\mathcal{K}(\mathcal{H})$.\\
\hlt{Key Insight}: Invertible modulo compact operators are closed under products.
\end{proof}
\end{corollary}

\begin{corollary} \label{cor:fredholm-index-zero-stable-compact-perturbations} \hlt{Compact Perturbations Preserve Index-Zero Fredholm Operators}\\
Let $\mathcal{H}$ be Hilbert and let $T,K\in\LL(\mathcal{H})$. If $T$ is Fredholm, $\operatorname{ind}T=0$, and $K$ is compact, then:
\begin{enumerate}[label=(\roman*)]
\setlength{\itemsep}{0pt}
\item $T+K$ is Fredholm.
\item $\operatorname{ind}(T+K)=0$.
\end{enumerate}
\begin{proof}
By Theorem \ref{thm:fredholm-index-zero-invertible-plus-compact}, $T=S+C$, $S$ invertible, $C$ compact. Set $\widetilde C\defeq C+K$.\\
By Proposition \ref{prop:compact-operators-linear-ideal}, $\widetilde C\in\mathcal{K}(\mathcal{H})$ and $T+K=S+\widetilde C$.
\begin{enumerate}[label=(\roman*)]
\setlength{\itemsep}{0pt}
\item $T+K=S+\widetilde C$, $S$ invertible, $\widetilde C$ compact $\Rightarrow T+K$ Fredholm by Theorem \ref{thm:fredholm-index-zero-invertible-plus-compact}.
\item $T+K=S+\widetilde C$, $S$ invertible, $\widetilde C$ compact $\Rightarrow\operatorname{ind}(T+K)=0$ by Theorem \ref{thm:fredholm-index-zero-invertible-plus-compact}.
\end{enumerate}
\hlt{Structural Commentary}: $T=S+C$, $K\in\mathcal{K}(\mathcal{H})\Rightarrow T+K=S+\widetilde C$, $\widetilde C=C+K\in\mathcal{K}(\mathcal{H})$.\\
\hlt{Key Insight}: $\operatorname{ind}=0$ is stable under compact perturbation.
\end{proof}
\end{corollary}
